\documentclass[11pt]{article}
\usepackage[margin=1in]{geometry}
\usepackage[T1]{fontenc}
\usepackage[utf8]{inputenc}
\usepackage{lmodern}
\usepackage{amsmath,amssymb,amsthm,mathtools}
\usepackage{booktabs,array,enumitem}
\usepackage{xcolor}
\usepackage[hidelinks]{hyperref}
\usepackage{microtype}

\newcommand{\Logic}{\mathcal{ALCI}_{\mathrm{RCC5}}}
\newcommand{\EQ}{\mathsf{EQ}}
\newcommand{\DR}{\mathsf{DR}}
\newcommand{\PO}{\mathsf{PO}}
\newcommand{\PP}{\mathsf{PP}}
\newcommand{\PPI}{\mathsf{PPI}}
\newcommand{\Base}{\mathsf{B}}
\newcommand{\comp}{\mathsf{comp}}
\newcommand{\Cl}{\mathsf{cl}}
\newcommand{\Mos}{\mathsf{Mos}}
\newcommand{\Q}{\mathcal Q}
\newcommand{\TQ}{\mathcal T_{\Q}}
\newcommand{\Pos}{\mathsf{Pos}}
\newcommand{\Trip}{\mathsf{Trip}}
\newcommand{\Side}{\mathsf{Side}}
\newcommand{\Reach}{\mathrel{\mathsf{Reach}}}
\newcommand{\reach}{\mathrel{\mathsf{reach}}}
\newcommand{\eqc}{\equiv_{\EQ}}
\newcommand{\blkeq}{\equiv_{\mathrm{blk}}}
\newcommand{\Stratum}{\mathsf{Stratum}}
\newcommand{\Subsrc}{\mathsf{Subsrc}}
\newcommand{\Vrt}{\mathsf{Vert}}
\newcommand{\Port}{\mathsf{Port}}
\newcommand{\lab}{\mathsf{lab}}
\newcommand{\vertical}{\mathsf{vert}}
\newcommand{\sibling}{\mathsf{sib}}
\newcommand{\Inv}{\mathsf{Inv}}
\newcommand{\U}{\mathcal U}
\newcommand{\md}{\mathsf{md}}

\title{Formal Review of the Round-6 Consolidated Generated Split-Forest Proof}
\author{Review prepared by GPT-5.5 Pro}
\date{May 27, 2026}

\begin{document}
\sloppy
\maketitle

\begin{abstract}
This memorandum reviews the attached manuscript
\texttt{repaired\_split\_forest\_no\_automata\_proof\_v2\_consolidated\_round6}
and its TeX source.  The round-6 document is a clear improvement over the
previous consolidated version.  It compiles cleanly, it directly addresses the
flat side-frontier problem by introducing recursive verticalization, it replaces
per-triple mosaic coverage by an explicit abstract label function $L_{\Q}$, it
adds cross-pair universal safety, and it strengthens the request-cycle argument
with half-open cycles and extended profiles.  Nevertheless, I would not yet
regard the decidability proof as closed.  The most serious remaining issue is
that the new abstract-position labelling mechanism appears to reintroduce the
blocking/equality conflation it was designed to avoid: if all occurrences with
the same abstract position symbol are labelled by $L_{\Q}(\pi,\pi)=\EQ$, then
distinct laps of a blocking cycle become semantically equal.  This affects the
basic one-state upward-chain certificate.  Additional remaining issues concern
the occurrence-local interpretation of equality ports, the precise stratum
identity stored in the finite certificate, the recursive side-interface
construction, universal safety across pumped cycles, and inconsistent side-width
bounds.  My verdict is therefore: substantial progress, but not yet a fully
watertight decidability proof.
\end{abstract}

\tableofcontents

\section{Materials reviewed}

I reviewed the following uploaded files:
\begin{itemize}[leftmargin=2em]
  \item \texttt{repaired\_split\_forest\_no\_automata\_proof\_v2\_consolidated\_round6.pdf};
  \item \texttt{repaired\_split\_forest\_no\_automata\_proof\_v2\_consolidated\_round6.tex}.
\end{itemize}
The TeX source compiles locally without fatal errors and produces a 39-page PDF.
I did not receive a new verification archive with this round, so this review is
a manuscript review rather than a rerun of an updated executable verification
suite.

\section{Executive verdict}

The round-6 manuscript makes real progress.  In particular, it no longer tries
to defend the old flat side-frontier discipline.  The new stratified side
interface and recursive verticalization discipline (M6$'$) are the right kind
of repair for side witnesses that are themselves $\PP/\PPI$-comparable.  The
manuscript also correctly identifies the previous inconsistency between
per-position support closure and per-triple mosaic coverage, and it attempts to
move the composition proof to the explicit global label function $L_{\Q}$.  The
request-cycle section is also more coherent than in the previous round.

However, several of the new repairs are not yet internally consistent.  The
most serious defect is independent of the support-mosaic details.  The current
combination of
\[
   q(u)\in\Pos_{\Q},\qquad
   \lab(u,v):=L_{\Q}(q(u),q(v)),\qquad
   L_{\Q}(\pi,\pi)=\EQ
\]
seems to force every two occurrences with the same finite abstract position
symbol to be $\EQ$-related.  Since blocking cycles necessarily reuse finite
abstract position symbols on fresh semantic occurrences, this collapses cyclic
blocking into semantic equality.  This directly conflicts with the manuscript's
central slogan that blocking is not equality.

Thus my revised verdict is:
\[
\boxed{\text{Round 6 is an important improvement, but the proof still does not stand as written.}}
\]
The proof architecture remains credible, but the formal certificate semantics
needs another occurrence-level repair.

\section{Positive changes in round 6}

\subsection{Stratified side interfaces address the old flat-frontier problem}

The earlier manuscript treated side/frontier positions too flatly.  That lost
satisfiable configurations in which two $\DR/\PO$ side witnesses of the same
source are themselves $\PP/\PPI$-comparable.  The round-6 text now explicitly
recognizes this and introduces a stratum forest inside a mosaic.  This is the
right conceptual repair: a pair of comparable side witnesses should be
co-vertical inside a nested local stratum rather than rejected as an illegal
sibling-frontier pair.

\subsection{The proof no longer depends on per-triple mosaic coverage}

The previous version weakened support closure to per-position support but later
continued to appeal to per-triple mosaic coverage.  The round-6 manuscript
tries to repair this by using the declared global pair-label function
\[
  L_{\Q}:\Pos_{\Q}^{2}\to\Base
\]
and enforcing the RCC5 composition table directly by validity clause (V11).
This is a promising direction.  If the abstract pair labels are defined over
the correct kind of occurrence-pair shapes, then composition can indeed be
checked finitely without requiring every triple to appear inside one mosaic.

\subsection{Cross-pair universal safety is now explicit}

The new validity clause (V15) is also the right sort of fix.  RCC5 composition
alone does not propagate DL types; if an overlap cross-label is declared to be
$R$, then universal restrictions $\forall R.D$ at the source must be checked
against the target type.  Making this a finite validity condition is an
improvement.

\subsection{The request-cycle section is more precise}

The switch to half-open cycle words $[i,j)$ and the use of extended profiles
including pending requests are both improvements.  This avoids one of the
previous defects, where a repeated ordinary profile might not carry the same
unfulfilled eventualities.

\section{Main remaining blocker: abstract positions collapse blocking into equality}
\label{sec:diagonal}

The central formal problem is the new definition of concrete labels through the
finite abstract position map.  The manuscript defines an abstract position
symbol $q(u)\in\Pos_{\Q}$ for each concrete occurrence $u$ and then defines
\[
  \lab(u,v):=L_{\Q}(q(u),q(v)).
\]
It also requires
\[
  L_{\Q}(\pi,\pi)=\EQ.
\]
This is harmless only if $q(u)=q(v)$ implies $u=v$.  But in a finite regular
certificate, this implication is false.  Blocking deliberately creates
infinitely many fresh semantic occurrences with the same finite profile, port,
context, and stratum data.

\subsection{Concrete failure on the upward-chain stress case}

Consider again
\[
  C_{\uparrow}=\exists\PP.\top\ \sqcap\ \forall\PP.\exists\PP.\top.
\]
The intended regular certificate is the one-state upward loop whose unfolding is
\[
  u_0\PP u_1\PP u_2\PP\cdots,
\]
with all $u_i$ fresh occurrences of the same finite state/profile.  In such a
certificate, the occurrences $u_i$ and $u_j$ for $i\ne j$ will have the same
abstract position symbol:
\[
  q(u_i)=q(u_j)=\pi.
\]
The round-6 labelling rule then gives
\[
  \lab(u_i,u_j)=L_{\Q}(\pi,
  \pi)=\EQ.
\]
This is impossible.  The intended relation for $i<j$ is $\PP$, and strong
$\EQ$ semantics forbids two distinct occurrences from being $\EQ$-related.
If the quotient identifies them, the strict $\PP$ chain collapses into a
$\PP$-cycle.  If the quotient does not identify them, the pre-quotient or
quotient frame has $\EQ$ between distinct elements.  Either way the current
formal definitions do not represent the very blocking cycle they are meant to
represent.

\subsection{Why this is not just notation}

This defect affects several later arguments:
\begin{itemize}[leftmargin=2em]
  \item the proof of the support-closed patchwork theorem uses
        $\lab(u,v)=L_{\Q}(q(u),q(v))$ for arbitrary concrete pairs;
  \item Lemma 7.4/JEPD relies on $\EQ$ iff same quotient class;
  \item the typed-$\EQ$ quotient proof assumes blocking repetition is not
        semantic equality;
  \item V9 forbids equality between strict vertical comparables.
\end{itemize}
These cannot all be true if $L_{\Q}(\pi,\pi)=\EQ$ is applied to two distinct
fresh occurrences with the same abstract symbol.

\subsection{Required repair}

The label function cannot be a plain function
\[
  L_{\Q}:\Pos_{\Q}^{2}\to\Base
\]
with diagonal $\EQ$ unless abstract positions are occurrence-unique, which would
make the certificate infinite.  A finite certificate needs a label function over
\emph{abstract occurrence-pair shapes}, not just over pairs of position symbols.
At minimum the pair shape must distinguish:
\begin{enumerate}[leftmargin=2em,label=(D\arabic*)]
  \item the literal identity pair $(u,u)$, where the label is $\EQ$;
  \item two distinct occurrences with the same abstract position symbol but in
        different laps or different context instances;
  \item vertical successor/predecessor copies across a blocking cycle;
  \item side/mosaic-visible pairs;
  \item explicit equality-port pairs.
\end{enumerate}
One possible replacement is a finite function
\[
  L_{\Q}^{\mathrm{pair}}:\mathsf{PairShape}_{\Q}\to\Base
\]
where $\mathsf{PairShape}_{\Q}$ contains an explicit identity flag and enough
relative-copy information to distinguish same-symbol identity from same-symbol
fresh copies.  The validity clause (V11) would then quantify over compatible
\emph{triple shapes}, not over $\Pos_{\Q}^{3}$ alone.

\section{Equality-port lifting remains ambiguous}
\label{sec:eqports}

A related issue occurs in the definition of the occurrence lift of
$\eqc^{Q}$.  The manuscript says that $\eqc^{Q}$ is the reflexive, symmetric,
transitive closure of explicit port links, and then writes informally that
\[
  u\eqc v \quad\text{iff}\quad \Port(u)\eqc^{Q}\Port(v)
\]
within the same generated context or transitively across explicit links.
If this is read literally at the state-port level, then every occurrence with
port $(s,p)$ is equivalent to every other occurrence with port $(s,p)$, because
$\eqc^{Q}$ is reflexive.  Again, distinct laps of a blocking cycle collapse.

There is also tension with V9.b.  The manuscript says that a finite cycle counts
as strict reachability because each lap unfolds to a fresh occurrence.  Thus a
self-loop can make
\[
  (s,p)\reach_{\PP}(s,p).
\]
But $\eqc^{Q}$ is reflexive on $(s,p)$, so V9.b, read literally at the finite
state-port level, would reject the very self-loop needed for
$C_{\uparrow}$.  If V9.b is not meant to reject the reflexive port occurrence,
then the text must say so and must define the occurrence-level lift explicitly.

\paragraph{Required repair.}
The finite port relation should be treated as a \emph{template} for explicit
local equality edges generated inside a context instance.  Its lift should be:
\[
  u\eqc v
  \quad\Longleftrightarrow\quad
  u=v\ \text{or there is a path of instantiated explicit equality edges between }u\text{ and }v.
\]
It should not identify all occurrences with the same finite port.  V9.b should
then be stated over instantiated occurrence-level equality edges, or over pair
shapes that can actually occur in one context instance, not over reflexive
state-port equality in the finite quotient.

\section{Stratum identity is still not fully present in the finite abstraction}

The abstract position definition still uses a stratum field
\[
  \zeta\in\{\vertical,\sibling\}.
\]
But M6$'$ is no longer a binary vertical/sibling distinction.  It uses a full
stratum forest
\[
  \mathcal V=(\Stratum,\Subsrc)
\]
with nested strata.  Whether two positions are co-vertical depends on the
identity of the stratum containing them, not merely on whether each is marked
``vertical''.  Two positions can both be vertical but in different strata; they
are then cross-stratum and should not receive $\PP/\PPI$ labels.

The introduction says the round-6 certificate grammar tracks the full stratum
forest object $\mathcal V$, but the formal abstract-position definition still
looks round-4-like.  The permitted-pair condition says that co-verticality is
``lifted via the certificate's stratum-membership data'', but that data is not
part of the abstract position symbol as written.

\paragraph{Required repair.}
Replace the binary field $\zeta\in\{\vertical,\sibling\}$ by a stratum identity
or stratum-path field, for example
\[
  \zeta \in \mathsf{StratumId}(C),
\]
where $C$ is the context scheme.  Then define co-verticality in
$\Pos_{\Q}$ by equality of an actual stratum identifier.  Validity clauses
(V2), (V11), (V12), and the enumeration grammar should all use the same
stratum-forest object, not the old $\Vrt$ notation or a binary vertical/sibling
flag.

\section{The recursive side-interface construction is improved but still underspecified}

The new side-interface legality lemma is much closer to the needed statement.
However, the construction of nested strata is still ambiguous for configurations
with more than two comparable side positions.

The lemma says that for every side witness $w$ whose original is $\PP/\PPI$-
related to another side position $w'$, a nested stratum $V_w$ is introduced.
But if one side position is comparable to several others, or if the side
positions form a chain
\[
  w_0\PP w_1\PP w_2,
\]
it is not clear whether there is one stratum per comparable pair, one stratum
per source $w$, or a recursively nested chain of strata.  The stratum forest
axioms allow overlaps only at parent-source positions, so the construction must
specify exactly how such chains are represented without creating illegal
stratum intersections.

There is also a tension inside M6$'$.  Clause (M6$'$.b) says that a containment
label in a stratum must match the source/mate versus selected-witness structure.
That appears to allow $\PP/\PPI$ labels only between the stratum source (or its
mates) and selected witnesses.  But the vertical-realisation proof later has a
subcase where both endpoints are selected containment witnesses of the same
stratum and $L(p,q)\in\{\PP,\PPI\}$.  As written, that subcase seems forbidden
by M6$'$.b rather than explained by it.

\paragraph{Required repair.}
Each stratum should carry an explicit local containment order or local parent
map, not merely a source and a set of selected witnesses.  Then M6$'$.b can say:
$L(p,q)\in\{\PP,\PPI\}$ iff $p,q$ are comparable in the local stratum order,
possibly modulo equality mates.  The recursive side-interface lemma should
construct this local order for arbitrary finite comparable side-witness
networks and prove the stated modal-depth width recurrence.

\section{Request cycles handle existentials better, but universal wrap-around still needs proof}

The request-cycle lemma now handles transitive existential requests more
carefully.  However, pumping a half-open segment can create new future
successors for a universal formula.  If a node $u_k$ in the cycle has
$\forall R.D$, then in the pumped path it sees not only later nodes
$u_{k+1},\ldots,u_{j-1}$ from the same lap, but also the next lap
$u_i,u_{i+1},\ldots$.  The original infinite path did not require $u_i,
\ldots,u_{k-1}$ to satisfy $D$, because they were earlier than $u_k$.

The proof currently says that universal truth is preserved because the visible
neighbourhood is the same as in the original path.  That is not literally true
under cyclic pumping: the wrap-around part of the next lap is newly visible as
a strict future successor.

This may be repairable.  The profile already gestures toward storing formulas
true at all/some strict ancestors and descendants.  But the request-cycle lemma
must use that information explicitly.  A valid cycle should satisfy a finite
universal-cycle safety condition:
\[
  \forall R.D\in\lambda(u_k)
  \quad\Longrightarrow\quad
  D\in\lambda(u_m)
  \quad\text{for every later position }m\text{ in the cyclic order.}
\]
For a half-open cycle $[i,j)$ this includes the wrap-around positions
$i\le m<k$ in the next lap.  The completeness proof then needs to show that a
cycle satisfying this stronger condition can always be selected from a
satisfying ray, perhaps by choosing the lasso after all universal-suffix
summaries have stabilized.

\section{The side-width and enumeration bounds are internally inconsistent}

The round-6 introduction and Lemma 4.5 say that the old cubic side-width bound
has been replaced by a modal-depth recurrence
\[
  m_*(C_0)=m(\md(C_0)).
\]
The enumeration lemma also says ``No cubic side-width is used.''  But the later
``Explicit component bounds'' section still defines
\[
  m(n)=1+n^3+3n^2+3n=O(n^3)
\]
as a cubic polynomial bound on $|\Side(u)|$ and on mosaic arity.  The same
section then defines polynomial bounds using this cubic $m(n)$.

This is not merely a cosmetic mismatch.  The finite search bound is the final
step of the decidability proof.  The proof needs one consistent bound, with the
same symbol used throughout.  If the correct bound is the modal-depth recurrence,
then the grammar and enumeration estimates should use $m_*(C_0)$ everywhere.
If a cubic bound is actually sufficient, then Lemma 4.5 and the surrounding
round-6 explanation should be changed accordingly.

\section{The Renz--Nebel / patchwork layer should be simplified or restated}

The manuscript defines semantics as strong abstract atomic RCC5 frames, not as
topological or set-theoretic region models.  For this abstract semantics,
soundness only needs JEPD, inverse symmetry, and RCC5 composition closure.  The
new $L_{\Q}$-direct proof already aims to establish exactly these abstract frame
conditions.

The Renz--Nebel realisability theorem is therefore not obviously needed for the
main decidability theorem as stated.  Worse, the ``Moreover'' clause in the
current restatement is too weak as written: saying there exists some
$\ell\in\comp(\rho_1(u,v),\rho_2(v,w))$ does not itself define a globally
consistent cross-label for all mixed triples.  The later strong-overlap
condition supplies actual cross-labels, so the theorem statement should not
suggest that mere non-emptiness of a composition entry is enough.

\paragraph{Recommendation.}
Either remove the set-theoretic Renz--Nebel realisation layer from the main
proof, or clearly mark it as optional background.  For the abstract-frame
result, the patchwork theorem can be replaced by the direct statement:
if the certificate supplies a well-defined pair-shape labelling satisfying JEPD,
inverse symmetry, and V11 on all compatible triple shapes, then the unfolded
structure is an abstract RCC5 frame.

\section{Recommended repairs before claiming decidability}

I recommend the following concrete repair sequence.

\begin{enumerate}[leftmargin=2em,label=(R\arabic*)]
  \item Replace $L_{\Q}:\Pos_{\Q}^{2}\to\Base$ by a label function on
        finite abstract occurrence-pair shapes.  Include an explicit identity
        flag so that $\EQ$ on the diagonal applies only to literal identity,
        not to two fresh blocked copies with the same finite position symbol.
  \item Redefine the occurrence lift of equality ports.  Equality links should
        be instantiated local edges inside generated context occurrences;
        reflexivity of a finite port should not identify all occurrences using
        that port.
  \item Rewrite V9.b at the occurrence-template level.  It should forbid
        explicit equality between two positions that can be strict-vertical
        comparable in the same unfolding instance, without rejecting self-loop
        blocking profiles by reflexivity.
  \item Make stratum identity part of abstract positions or pair shapes.  A
        binary vertical/sibling flag is not enough for recursive M6$'$.
  \item Strengthen M6$'$ so each stratum has an explicit local containment
        order, and rewrite the side-interface construction for arbitrary finite
        comparable side-witness networks.
  \item Add a universal-cycle safety lemma for pumped half-open cycles, or
        strengthen the extended profile so that request-fulfilled cycles are
        also universal-safe across wrap-around.
  \item Normalize the side-width bound: use either the modal-depth recurrence
        $m_*(C_0)$ everywhere or prove the cubic bound everywhere.
  \item Simplify the patchwork layer to the abstract-frame needs of the theorem,
        or restate the Renz--Nebel component precisely and mark which parts are
        optional.
\end{enumerate}

\section{Conclusion}

Round 6 is genuine progress.  The manuscript has moved from the old flat
side-interface picture toward the right recursive, occurrence-based split-forest
architecture.  The major conceptual repairs are in place: witness-generated
covers, split copies, typed equality, recursive verticalization, cross-pair
universal safety, and request-fulfilled cycles.

But the proof still contains a serious formal defect: the finite abstract
position labelling currently equates distinct blocked copies by assigning
$L_{\Q}(\pi,\pi)=\EQ$ whenever their abstract position symbols coincide.  This
undermines the central blocking-not-equality principle and breaks the canonical
one-state certificate for
\[
  \exists\PP.\top\sqcap\forall\PP.\exists\PP.\top.
\]
Until the label function and equality-port lift are repaired at the
occurrence-pair-shape level, I would not regard the decidability theorem as
established.

The architecture is still promising, but the next manuscript should focus on
this single foundational correction before further polishing the mosaic and
verification layers.

\end{document}
